\begin{thebibliography}{99}
\setlength{\itemsep}{2pt}

\bibitem[Agarwal et~al.(2025)]{agarwal2025} Agarwal V, Inoue F, Schubach M, et al. (2025). Massively parallel characterization of transcriptional regulatory elements. \textit{Nature} 639:411–420. \href{https://doi.org/10.1038/s41586-024-08430-9}{doi:10.1038/s41586-024-08430-9}

\bibitem[Ashuach et~al.(2019)]{ashuach2019} Ashuach T, Fischer DS, Kreimer A, et al. (2019). MPRAnalyze: statistical framework for massively parallel reporter assays. \textit{Genome Biology} 20:183. \href{https://doi.org/10.1186/s13059-019-1787-z}{doi:10.1186/s13059-019-1787-z}

\bibitem[Avsec et~al.(2021)]{avsec2021} Avsec Ž, Agarwal V, Visentin D, et al. (2021). Effective gene expression prediction from sequence by integrating long-range interactions. \textit{Nature Methods} 18:1196–1203. \href{https://doi.org/10.1038/s41592-021-01252-x}{doi:10.1038/s41592-021-01252-x}

\bibitem[Aygün et~al.(2021)]{aygun2021} Aygün N, Elwell AL, Liang D, et al. (2021). Brain-trait-associated variants impact cell-type-specific gene regulation during neurogenesis. \textit{American Journal of Human Genetics} 108:1647–1668. \href{https://doi.org/10.1016/j.ajhg.2021.07.011}{doi:10.1016/j.ajhg.2021.07.011}

\bibitem[Balaskas et~al.(2012)]{balaskas2012} Balaskas N, Ribeiro A, Panovska J, et al. (2012). Gene regulatory logic for reading the Sonic Hedgehog signaling gradient in the vertebrate neural tube. \textit{Cell} 148:273–284. \href{https://doi.org/10.1016/j.cell.2011.10.047}{doi:10.1016/j.cell.2011.10.047}

\bibitem[Beal et~al.(2011)]{beal2011} Beal J, Lu T, Weiss R. (2011). Automatic compilation from high-level biologically-oriented programming language to genetic regulatory networks. \textit{PLoS ONE} 6:e22490. \href{https://doi.org/10.1371/journal.pone.0022490}{doi:10.1371/journal.pone.0022490}

\bibitem[Bergmann et~al.(2014)]{bergmann2014} Bergmann FT, Adams R, Moodie S, et al. (2014). COMBINE archive and OMEX format: one file to share all information to reproduce a modeling project. \textit{BMC Bioinformatics} 15:369. \href{https://doi.org/10.1186/s12859-014-0369-z}{doi:10.1186/s12859-014-0369-z}

\bibitem[Bilitchenko et~al.(2011)]{bilitchenko2011} Bilitchenko L, Liu A, Cheung S, et al. (2011). Eugene – a domain specific language for specifying and constraining synthetic biological parts, devices, and systems. \textit{PLoS ONE} 6:e18882. \href{https://doi.org/10.1371/journal.pone.0018882}{doi:10.1371/journal.pone.0018882}

\bibitem[Cai et~al.(2009)]{cai2009} Cai Y, Lux MW, Adam L, Peccoud J. (2009). Modeling structure-function relationships in synthetic DNA sequences using attribute grammars. \textit{PLoS Computational Biology} 5:e1000529. \href{https://doi.org/10.1371/journal.pcbi.1000529}{doi:10.1371/journal.pcbi.1000529}

\bibitem[Chen et~al.(2022)]{chen2022} Chen KM, Wong AK, Troyanskaya OG, Zhou J. (2022). A sequence-based global map of regulatory activity for deciphering human genetics. \textit{Nature Genetics} 54:940–949. \href{https://doi.org/10.1038/s41588-022-01102-2}{doi:10.1038/s41588-022-01102-2}

\bibitem[Cohen et~al.(2014)]{cohen2014} Cohen M, Page KM, Perez-Carrasco R, Barnes CP, Briscoe J. (2014). A theoretical framework for the regulation of Shh morphogen-controlled gene expression. \textit{Development} 141:3868–3878. \href{https://doi.org/10.1242/dev.112573}{doi:10.1242/dev.112573}

\bibitem[Deng et~al.(2024)]{deng2024} Deng C, Whalen S, Steyert M, et al. (2024). Massively parallel characterization of regulatory elements in the developing human cortex. \textit{Science} 384:eadh0559. \href{https://doi.org/10.1126/science.adh0559}{doi:10.1126/science.adh0559}

\bibitem[Frankish et~al.(2021)]{frankish2021} Frankish A, et al. (2021). GENCODE 2021. \textit{Nucleic Acids Research} 49:D916–D923. \href{https://doi.org/10.1093/nar/gkaa1087}{doi:10.1093/nar/gkaa1087}

\bibitem[Fromer et~al.(2016)]{fromer2016} Fromer M, Roussos P, Sieberts SK, et al. (2016). Gene expression elucidates functional impact of polygenic risk for schizophrenia. \textit{Nature Neuroscience} 19:1442–1453. \href{https://doi.org/10.1038/nn.4399}{doi:10.1038/nn.4399}

\bibitem[Gneiting and Raftery(2007)]{gneiting2007} Gneiting T, Raftery AE. (2007). Strictly proper scoring rules, prediction, and estimation. \textit{Journal of the American Statistical Association} 102:359–378. \href{https://doi.org/10.1198/016214506000001437}{doi:10.1198/016214506000001437}

\bibitem[Gong et~al.(2020)]{gong2020} Gong et al. (2020). The requirement of Sox2 for the spinal cord motor neuron development of zebrafish. \textit{Frontiers in Molecular Neuroscience} 13:34. \href{https://doi.org/10.3389/fnmol.2020.00034}{doi:10.3389/fnmol.2020.00034}

\bibitem[GTEx Consortium(2020)]{gtex2020} GTEx Consortium (2020). The GTEx Consortium atlas of genetic regulatory effects across human tissues. \textit{Science} 369:1318–1330. \href{https://doi.org/10.1126/science.aaz1776}{doi:10.1126/science.aaz1776}

\bibitem[Guo et~al.(2023)]{guo2023} Guo et al. (2023). Integrative analyses highlight functional regulatory variants associated with neuropsychiatric diseases. \textit{Nature Genetics} 55:1876–1891. \href{https://doi.org/10.1038/s41588-023-01533-5}{doi:10.1038/s41588-023-01533-5}

\bibitem[Hallikas et~al.(2006)]{hallikas2006} Hallikas O, Palin K, Sinjushina N, et al. (2006). Genome-wide prediction of mammalian enhancers based on analysis of transcription-factor binding affinity. \textit{Cell} 124:47–59. \href{https://doi.org/10.1016/j.cell.2005.10.042}{doi:10.1016/j.cell.2005.10.042}

\bibitem[Herrera-Delgado et~al.(2020)]{herreradelgado2020} Herrera-Delgado E, Briscoe J, Sollich P. (2020). Tractable nonlinear memory functions as a tool to capture and explain dynamical behaviors. \textit{Physical Review Research} 2:043069. \href{https://doi.org/10.1103/PhysRevResearch.2.043069}{doi:10.1103/PhysRevResearch.2.043069}

\bibitem[Huber(1964)]{huber1964} Huber PJ. (1964). Robust estimation of a location parameter. \textit{Annals of Mathematical Statistics} 35:73–101. \href{https://doi.org/10.1214/aoms/1177703732}{doi:10.1214/aoms/1177703732}

\bibitem[Jaffe et~al.(2018)]{jaffe2018} Jaffe AE, Straub RE, Shin JH, et al. (2018). Developmental and genetic regulation of the human cortex transcriptome illuminate schizophrenia pathogenesis. \textit{Nature Neuroscience} 21:1117–1125. \href{https://doi.org/10.1038/s41593-018-0197-y}{doi:10.1038/s41593-018-0197-y}

\bibitem[Jerber et~al.(2021)]{jerber2021} Jerber J, Seaton DD, Cuomo ASE, et al. (2021). Population-scale single-cell RNA-seq profiling across dopaminergic neuron differentiation. \textit{Nature Genetics} 53:304–312. \href{https://doi.org/10.1038/s41588-021-00801-6}{doi:10.1038/s41588-021-00801-6}

\bibitem[Kang et~al.(2011)]{kang2011} Kang HJ, Kawasawa YI, Cheng F, et al. (2011). Spatio-temporal transcriptome of the human brain. \textit{Nature} 478:483–489. \href{https://doi.org/10.1038/nature10523}{doi:10.1038/nature10523}

\bibitem[Kelley et~al.(2016)]{kelley2016} Kelley DR, Snoek J, Rinn JL. (2016). Basset: learning the regulatory code of the accessible genome with deep convolutional neural networks. \textit{Genome Research} 26:990–999. \href{https://doi.org/10.1101/gr.200535.115}{doi:10.1101/gr.200535.115}

\bibitem[Kelley et~al.(2018)]{kelley2018} Kelley DR, Reshef YA, Bileschi M, et al. (2018). Sequential regulatory activity prediction across chromosomes with convolutional neural networks. \textit{Genome Research} 28:739–750. \href{https://doi.org/10.1101/gr.227819.117}{doi:10.1101/gr.227819.117}

\bibitem[Kerimov et~al.(2021)]{kerimov2021} Kerimov N, Hayhurst JD, Peikova K, et al. (2021). A compendium of uniformly processed human gene expression and splicing quantitative trait loci. \textit{Nature Genetics} 53:1290–1299. \href{https://doi.org/10.1038/s41588-021-00924-w}{doi:10.1038/s41588-021-00924-w}

\bibitem[Kircher et~al.(2019)]{kircher2019} Kircher M, Xiong C, Martin B, et al. (2019). Saturation mutagenesis of twenty disease-associated regulatory elements at single base-pair resolution. \textit{Nature Communications} 10:3583. \href{https://doi.org/10.1038/s41467-019-11526-w}{doi:10.1038/s41467-019-11526-w}

\bibitem[Kreimer et~al.(2022)]{kreimer2022} Kreimer A, Ashuach T, Inoue F, et al. (2022). Massively parallel reporter perturbation assays uncover temporal regulatory architecture during neural differentiation. \textit{Nature Communications} 13:1504. \href{https://doi.org/10.1038/s41467-022-28659-0}{doi:10.1038/s41467-022-28659-0}

\bibitem[Lee et~al.(2025)]{lee2025} Lee S, McAfee JC, Lee J, et al. (2025). Massively parallel reporter assay investigates shared genetic variants of eight psychiatric disorders. \textit{Cell} 188:1409–1424.e21. \href{https://doi.org/10.1016/j.cell.2024.12.022}{doi:10.1016/j.cell.2024.12.022}

\bibitem[Majoros et~al.(2020)]{majoros2020} Majoros WH, Kim YS, Barrera A, et al. (2020). Bayesian estimation of genetic regulatory effects in high-throughput reporter assays. \textit{Bioinformatics} 36:331–338. \href{https://doi.org/10.1093/bioinformatics/btz545}{doi:10.1093/bioinformatics/btz545}

\bibitem[Matoba et~al.(2026)]{matoba2026} Matoba N, McAfee JC, Krupa O, et al. (2026). Massively parallel assessment of gene regulatory activity at human cortical-structure-associated variants. \textit{Nature Neuroscience}. \href{https://doi.org/10.1038/s41593-026-02454-2}{doi:10.1038/s41593-026-02454-2}

\bibitem[McAfee et~al.(2023)]{mcafee2023} McAfee JC, Lee S, Lee J, et al. (2023). Systematic investigation of allelic regulatory activity of schizophrenia-associated common variants. \textit{Cell Genomics} 3:100404. \href{https://doi.org/10.1016/j.xgen.2023.100404}{doi:10.1016/j.xgen.2023.100404}

\bibitem[Melnikov et~al.(2012)]{melnikov2012} Melnikov A, Murugan A, Zhang X, et al. (2012). Systematic dissection and optimization of inducible enhancers in human cells using a massively parallel reporter assay. \textit{Nature Biotechnology} 30:271–277. \href{https://doi.org/10.1038/nbt.2137}{doi:10.1038/nbt.2137}

\bibitem[Miller et~al.(2014)]{miller2014} Miller JA, Ding SL, Sunkin SM, et al. (2014). Transcriptional landscape of the prenatal human brain. \textit{Nature} 508:199–206. \href{https://doi.org/10.1038/nature13185}{doi:10.1038/nature13185}

\bibitem[Myint et~al.(2019)]{myint2019} Myint L, Avramopoulos DG, Goff LA, Hansen KD. (2019). Linear models enable powerful differential activity analysis in massively parallel reporter assays. \textit{BMC Genomics} 20:209. \href{https://doi.org/10.1186/s12864-019-5556-x}{doi:10.1186/s12864-019-5556-x}

\bibitem[Ng et~al.(2017)]{ng2017} Ng B, White CC, Klein HU, et al. (2017). An xQTL map integrates the genetic architecture of the human brain's transcriptome and epigenome. \textit{Nature Neuroscience} 20:1418–1426. \href{https://doi.org/10.1038/nn.4632}{doi:10.1038/nn.4632}

\bibitem[Liu et~al.(2025)]{liu2025} Liu et al. (2025). Predicting the DNA binding specificity of transcription factor mutants using family-level biophysically interpretable machine learning. \textit{Nucleic Acids Research} 53:gkaf831. \href{https://doi.org/10.1093/nar/gkaf831}{doi:10.1093/nar/gkaf831}

\bibitem[Nielsen et~al.(2016)]{nielsen2016} Nielsen AAK, Der BS, Shin J, et al. (2016). Genetic circuit design automation. \textit{Science} 352:aac7341. \href{https://doi.org/10.1126/science.aac7341}{doi:10.1126/science.aac7341}

\bibitem[NIST(2015)]{nist2015} NIST (2015). \textit{Secure Hash Standard (SHS)}, FIPS 180-4. \url{https://csrc.nist.gov/pubs/fips/180-4/upd1/final}

\bibitem[Peterson et~al.(2012)]{peterson2012} Peterson KA, Nishi Y, Ma W, et al. (2012). Neural-specific Sox2 input and differential Gli-binding affinity provide context and positional information in Shh-directed neural patterning. \textit{Genes \& Development} 26:2802–2816. \href{https://doi.org/10.1101/gad.207142.112}{doi:10.1101/gad.207142.112}

\bibitem[Pollard et~al.(2010)]{pollard2010} Pollard KS, Hubisz MJ, Rosenbloom KR, Siepel A. (2010). Detection of nonneutral substitution rates on mammalian phylogenies. \textit{Genome Research} 20:110–121. \href{https://doi.org/10.1101/gr.097857.109}{doi:10.1101/gr.097857.109}

\bibitem[Ramasamy et~al.(2014)]{ramasamy2014} Ramasamy A, Trabzuni D, Guelfi S, et al. (2014). Genetic variability in the regulation of gene expression in ten regions of the human brain. \textit{Nature Neuroscience} 17:1418–1428. \href{https://doi.org/10.1038/nn.3801}{doi:10.1038/nn.3801}

\bibitem[Rauluseviciute et~al.(2024)]{rauluseviciute2024} Rauluseviciute I, Riudavets-Puig R, Blanc-Mathieu R, et al. (2024). JASPAR 2024: 20th anniversary of the open-access database of transcription factor binding profiles. \textit{Nucleic Acids Research} 52:D174–D182. \href{https://doi.org/10.1093/nar/gkad1059}{doi:10.1093/nar/gkad1059}

\bibitem[Roadmap Epigenomics Consortium(2015)]{roadmap2015} Roadmap Epigenomics Consortium, Kundaje A, Meuleman W, et al. (2015). Integrative analysis of 111 reference human epigenomes. \textit{Nature} 518:317–330. \href{https://doi.org/10.1038/nature14248}{doi:10.1038/nature14248}

\bibitem[Rummel et~al.(2023)]{rummel2023} Rummel CK, Gagliardi M, Ahmad R, et al. (2023). Massively parallel functional dissection of schizophrenia-associated noncoding genetic variants. \textit{Cell} 186:5165–5182.e33. \href{https://doi.org/10.1016/j.cell.2023.09.015}{doi:10.1016/j.cell.2023.09.015}

\bibitem[Rundgren et~al.(2020)]{rundgren2020} Rundgren A, Jordan B, Erdtman S. (2020). \textit{JSON Canonicalization Scheme (JCS)}, RFC 8785. \url{https://www.rfc-editor.org/rfc/rfc8785}

\bibitem[Salomon et~al.(2026)]{salomon2026} Salomon K, Deng C, Dash PM, et al. (2026). Massively parallel characterization and predictive modelling of neuronal regulatory variation. \textit{bioRxiv} (preprint, posted 17 July 2026). \href{https://doi.org/10.64898/2026.07.16.738760}{doi:10.64898/2026.07.16.738760}

\bibitem[Schwartzentruber et~al.(2018)]{schwartzentruber2018} Schwartzentruber J, Foskolou S, Kilpinen H, et al. (2018). Molecular and functional variation in iPSC-derived sensory neurons. \textit{Nature Genetics} 50:54–61. \href{https://doi.org/10.1038/s41588-017-0005-8}{doi:10.1038/s41588-017-0005-8}

\bibitem[Seitz et~al.(2024)]{seitz2024} Seitz EE, McCandlish DM, Kinney JB, Koo PK. (2024). Interpreting cis-regulatory mechanisms from genomic deep neural networks using surrogate models. \textit{Nature Machine Intelligence} 6:701–713. \href{https://doi.org/10.1038/s42256-024-00851-5}{doi:10.1038/s42256-024-00851-5}

\bibitem[Siraj et~al.(2026)]{siraj2026} Siraj L, et al. (2026). Functional dissection of complex trait variants at single-nucleotide resolution. \textit{Nature}. \href{https://doi.org/10.1038/s41586-026-10121-6}{doi:10.1038/s41586-026-10121-6}

\bibitem[Tareen et~al.(2022)]{tareen2022} Tareen A, Kooshkbaghi M, Posfai A, et al. (2022). MAVE-NN: learning genotype-phenotype maps from multiplex assays of variant effect. \textit{Genome Biology} 23:98. \href{https://doi.org/10.1186/s13059-022-02661-7}{doi:10.1186/s13059-022-02661-7}

\bibitem[Tewhey et~al.(2016)]{tewhey2016} Tewhey R, Kotliar D, Park DS, et al. (2016). Direct identification of hundreds of expression-modulating variants using a multiplexed reporter assay. \textit{Cell} 165:1519–1529. \href{https://doi.org/10.1016/j.cell.2016.04.027}{doi:10.1016/j.cell.2016.04.027}

\bibitem[Traets et~al.(2021)]{traets2021} Traets JJ, van der Burght SN, Rademakers S, et al. (2021). Mechanism of life-long maintenance of neuron identity despite molecular fluctuations. \textit{eLife} 10:e66955. \href{https://doi.org/10.7554/eLife.66955}{doi:10.7554/eLife.66955}

\bibitem[Vaishnav et~al.(2022)]{vaishnav2022} Vaishnav ED, de Boer CG, Molinet J, et al. (2022). The evolution, evolvability and engineering of gene regulatory DNA. \textit{Nature} 603:455–463. \href{https://doi.org/10.1038/s41586-022-04506-6}{doi:10.1038/s41586-022-04506-6}

\bibitem[Walker et~al.(2019)]{walker2019} Walker RL, Ramaswami G, Hartl C, et al. (2019). Genetic control of expression and splicing in developing human brain informs disease mechanisms. \textit{Cell} 179:750–771.e22. \href{https://doi.org/10.1016/j.cell.2019.09.021}{doi:10.1016/j.cell.2019.09.021}

\bibitem[Wang et~al.(2020)]{wang2020} Wang G, Sarkar A, Carbonetto P, Stephens M. (2020). A simple new approach to variable selection in regression, with application to genetic fine mapping. \textit{Journal of the Royal Statistical Society Series B} 82:1273–1300. \href{https://doi.org/10.1111/rssb.12388}{doi:10.1111/rssb.12388}

\bibitem[Weiss et~al.(2021)]{weiss2021} Weiss CV, Harshman L, Inoue F, et al. (2021). The cis-regulatory effects of modern human-specific variants. \textit{eLife} 10:e63713. \href{https://doi.org/10.7554/eLife.63713}{doi:10.7554/eLife.63713}

\bibitem[Wilson et~al.(2016)]{wilson2016} Wilson ML, et al. (2016). Genotype Specification Language. \textit{ACS Synthetic Biology} 5:471–478. \href{https://doi.org/10.1021/acssynbio.5b00194}{doi:10.1021/acssynbio.5b00194}

\bibitem[Young et~al.(2021)]{young2021} Young AMH, Kumasaka N, Calvert F, et al. (2021). A map of transcriptional heterogeneity and regulatory variation in human microglia. \textit{Nature Genetics} 53:861–868. \href{https://doi.org/10.1038/s41588-021-00875-2}{doi:10.1038/s41588-021-00875-2}

\bibitem[Zhou and Troyanskaya(2015)]{zhou2015} Zhou J, Troyanskaya OG. (2015). Predicting effects of noncoding variants with deep learning–based sequence model. \textit{Nature Methods} 12:931–934. \href{https://doi.org/10.1038/nmeth.3547}{doi:10.1038/nmeth.3547}

\bibitem[Zubler and Douglas(2009)]{zubler2009} Zubler F, Douglas R. (2009). A framework for modeling the growth and development of neurons and networks. \textit{Frontiers in Computational Neuroscience} 3:25. \href{https://doi.org/10.3389/neuro.10.025.2009}{doi:10.3389/neuro.10.025.2009}

\bibitem[Zubler et~al.(2011)]{zubler2011} Zubler F, Hauri A, Pfister S, et al. (2011). An instruction language for self-construction in the context of neural networks. \textit{Frontiers in Computational Neuroscience} 5:57. \href{https://doi.org/10.3389/fncom.2011.00057}{doi:10.3389/fncom.2011.00057}

\bibitem[Zubler et~al.(2013)]{zubler2013} Zubler F, Hauri A, Pfister S, et al. (2013). Simulating cortical development as a self constructing process: a novel multi-scale approach combining molecular and physical aspects. \textit{PLoS Computational Biology} 9:e1003173. \href{https://doi.org/10.1371/journal.pcbi.1003173}{doi:10.1371/journal.pcbi.1003173}

\end{thebibliography}
